\chapter{Excision Biopsy of a Breast Lump}
\index{Excision Biopsy of a Breast Lump}

\section*{Definition and Overview}
An excision biopsy is a surgical procedure in which an entire breast lump is removed so it can be examined under a microscope. It is performed when a lump needs a definitive diagnosis, most often when a core needle biopsy has been inconclusive, when the lump is small and entirely removable, or when a particular abnormality needs complete removal. The procedure is usually done under general anaesthetic as a day procedure, and most lumps turn out to be benign. If the lump is cancer, the excision biopsy is often the first step in treatment.

\section*{Epidemiology}
Breast lumps are very common, and the vast majority are benign. Many are investigated with imaging and core needle biopsy, which usually provides the diagnosis. Excision biopsy is needed in a smaller proportion of cases: when the needle biopsy is not definitive, when the imaging and biopsy do not agree, or when a lesion such as a papilloma is best removed entirely. Breast cancer is common, and accurate diagnosis of every lump is essential.

\section*{Aetiology and Pathophysiology}
Breast lumps have many causes.
\begin{itemize}
    \item \textbf{Benign causes:} fibroadenomas, cysts, fibrocystic change, and breast tissue density
    \item \textbf{Precancerous changes:} atypical hyperplasia and other changes that increase risk
    \item \textbf{Cancer:} breast cancer, which is why every persistent lump is investigated
    \item \textbf{Why excision:} removal of the whole lump allows the pathologist to examine all of it, giving a definitive diagnosis
    \item \textbf{Why not always:} most lumps are diagnosed adequately with imaging and needle biopsy, avoiding surgery
    \item \textbf{The procedure:} the lump is removed through a small incision, with the margins checked
\end{itemize}

\section*{Clinical Features}
Excision biopsy is recommended in specific situations.
\begin{itemize}
    \item A lump with an inconclusive needle biopsy
    \item A lump where the biopsy and imaging do not agree
    \item A small lump that is entirely removable
    \item A lesion that is best removed, such as a papilloma or atypical change
    \item An abnormality that grows or changes
    \item A suspicious lump in a person where imaging is difficult
\end{itemize}

\section*{Differential Diagnosis}
Before surgery, the lump is fully assessed.
\begin{itemize}
    \item Fibroadenoma and other benign lumps
    \item Breast cysts, which may be drained rather than excised
    \item Fibrocystic change
    \item Ductal carcinoma in situ and invasive breast cancer
    \item Infection or abscess
    \item Normal breast tissue or scar tissue
\end{itemize}

\section*{When to Seek Medical Care}
Any new breast lump should be assessed by a GP and referred to a breast clinic. Urgent care is needed for a lump that is growing rapidly, associated with skin changes, or accompanied by bloody nipple discharge.

\textbf{Contact your local emergency services for:}
\begin{itemize}
    \item Severe pain, swelling, or fever after the procedure, which may indicate infection
    \item Heavy bleeding from the wound
    \item Sudden difficulty breathing or chest pain after surgery
\end{itemize}

\section*{Diagnosis and Investigations}
Before the excision, the diagnosis is pursued with imaging and biopsy.
\begin{itemize}
    \item \textbf{Mammogram and ultrasound:} the first-line imaging
    \item \textbf{Core needle biopsy:} usually attempted first to obtain a diagnosis
    \item \textbf{Discussion at the breast clinic:} the decision to excise is made by a multidisciplinary team
    \item \textbf{Wire or marker localisation:} for lumps that cannot be felt, a wire or marker is placed to guide the surgery
    \item \textbf{Pre-operative assessment:} general health and anaesthetic review
\end{itemize}

\section*{Management}
The procedure and its aftermath are straightforward.
\begin{itemize}
    \item \textbf{Anaesthesia:} usually general anaesthetic, sometimes local
    \item \textbf{The operation:} a small incision over the lump, removal of the whole lump, and closure with stitches
    \item \textbf{Localisation:} for non-palpable lumps, a wire or marker is used
    \item \textbf{Pathology:} the lump is examined over the following week or two
    \item \textbf{Recovery:} day surgery with simple pain relief; return to normal activity within days
    \item \textbf{Results and follow-up:} the results are discussed, and further treatment is arranged if needed
\end{itemize}

\section*{Complications}
\begin{itemize}
    \item Bruising, swelling, and pain at the site
    \item Infection of the wound
    \item Bleeding and haematoma
    \item Scarring and changes in the shape of the breast
    \item Numbness around the scar
    \item If the lump is cancer, the possibility that further surgery is needed
\end{itemize}

\section*{Prognosis and Outcomes}
Excision biopsy provides a definitive diagnosis and is often curative for benign and precancerous conditions. If the lump is cancer, the excision biopsy has usually removed the entire lesion and forms the basis of further treatment planning. The recovery is quick, and the procedure is safe.

\section*{Prevention and Self-Care}
\begin{itemize}
    \item Be breast aware and report any new lump promptly
    \item Attend breast screening when invited
    \item Follow the post-operative wound care advice
    \item Keep the wound dry and clean
    \item Take simple pain relief as advised
    \item Attend the follow-up appointment to receive the results
    \item Seek care promptly for any signs of infection
\end{itemize}

\section*{Sources and Further Reading}
The following reputable sources provide further information for healthcare students and patients seeking reliable, current advice about excision biopsy of a breast lump.
\begin{itemize}
    \item BreastScreen Australia. \textit{Breast screening information.} https://www.health.gov.au/
    \item Cancer Council Australia. \textit{Breast cancer diagnosis and tests.} https://www.cancer.org.au/
    \item Healthdirect Australia. \textit{Breast lumps.} https://www.healthdirect.gov.au/breast-lumps
    \item NHS. \textit{Breast lump biopsy.} https://www.nhs.uk/conditions/breast-lump/
    \item Mayo Clinic. \textit{Breast biopsy.} https://www.mayoclinic.org/tests-procedures/breast-biopsy/
    \item Breast Cancer Network Australia. \textit{Information and support.} https://www.bcna.org.au/
\end{itemize}
